\speciestitle{Ozone}{O\cs3}{O3}%
\begin{description}
\item[Useful range:] 261\,--\,0.02\,hPa  
\item[Contact:] Lucien Froidevaux (stratosphere/mesosphere), 
  \email{Lucien.Froidevaux@jpl.nasa.gov} \\ 
  Michael Schwartz (upper troposphere), 
  \email{Michael.J.Schwartz@jpl.nasa.gov} 
\end{description} 
\label{sec:StandardO3} 
 
\subsection*{Introduction} 
 
The \thisv standard O\cs3 product is taken from the 240-GHz retrieval,
which provides the highest sensitivity down into the upper troposphere,
as well as in the mesosphere.  Table~\ref{tab:O3Summary} summarizes the
typical resolution, precision, and systematic uncertainty estimates as a
function of pressure.  Papers describing detailed validation of the MLS
v2.2 product and comparisons with other data sets were published in a
special Aura validation issue of the {\it Journal of Geophysical
  Research}, see \citep{FroidevauxEtal08O3,JiangYEtal07,LiveseyEtal08}.  In
the stratosphere and above, \thisv ozone profiles are very similar to the
v2.2 profiles, so the stratospheric results from the above references
will generally hold for the \thisv product.  Initial documentation of
changes, improvements, and issues with \thisv data are discussed here,
including data screening criteria (which are now somewhat more complex
than in v2.2 data).  The morphology of the (zonal mean) \thisv ozone data
appears to be more reliable (realistic) for the largest pressure values
where ozone is retrieved in the upper troposphere, notably at 316\,hPa
and also at 261\,hPa (a new level for \thisv). The 316\,hPa ozone
retrievals will, however, require further validation before we can
recommend this pressure surface for scientific use.
 
There are 2 separate stratospheric ozone columns (typically in very good
agreement) in the \texttt{L2GP-O3} files, with swath names
%
`\verb*+O3 column-MLS+'
%
and
%
`\verb*+O3 column-GEOS5+',
%
corresponding to the use of tropopause pressures (WMO definition)
determined from MLS or GEOS-5 temperatures, respectively.  Data users
can also provide their own calculations of column ozone values (with
better screening), based on the MLS ozone profiles, given that poorly
defined tropopause values can lead to relatively large scatter at
certain places (and times) for ozone column results.
 
\subsection*{Comparison of \thisv with v2.2} 
 
Between 316\,hPa and 1\,hPa, \thisv ozone profiles are retrieved on 12
surfaces per decade, a grid twice as fine as the 6-level-per-decade grid
used in v2.2. This finer grid makes possible some improvement in
vertical resolution in the UTLS, with retrieved precisions similar to
(or slightly larger than) the v2.2 values, but at the cost of poorer
horizontal resolution (see Table\ref{tab:O3Summary}.)  For pressures
less than 1\,hPa, the retrieval grid has not changed from the v2.2 grid;
it becomes coarser, with 6 surfaces per decade change in pressure
between 1 and 0.1\,hPa, and 3 surfaces per decade change in pressure
from 0.1\,hPa to 0.01\,hPa.  Other changes in the retrievals for the
240-GHz phase (affecting CO and HNO\cs3 as well as ozone) include the
manner in which the spectral baseline is modeled.  The \thisv retrieval
fits a relative-humidity-like, frequency squared extinction baseline at
the lowest retrieval levels, rather than a spectrally flat extinction
profile, as used in v2.2.  This modification reduces ozone biases in the
upper troposphere in moist clear air (or thin cloud) conditions, but
gives a poorer fit to the impacts of thick clouds (scattering from ice
particles in convective cores).  Cloud effects lead to more scatter and
vertically-oscillating profiles in \thisv, for the most part in the
tropics, and methods to screen out cloud-impacted profiles are discussed
below.  Also, there have been relatively minor spectroscopic changes for
these \thisv ozone retrievals: the line mixing parameters are now taken
from the recent (unpublished) work of {\it DeLucia et al.}.
 
\begin{figure} 
\centering\dontincludegraphics[angle=90,width=\linewidth]%
{figures/InspectZMContour_O3--v2-2_O3--v3-3_2006m11_LatPress-All-NaN-Allheights} 
\caption{Zonal averages for stratospheric and mesospheric MLS ozone
  profiles during November, 2006, showing the MLS v2.2 ozone mixing
  ratio contours (top left panel), the \thisv contours (top right panel),
  and their differences in ppmv (\thisv minus v2.2, bottom left panel) and
  percent (\thisv minus v2.2 versus v2.2, bottom right panel). }
\label{fig:MLSv2v3O3ContourStrat} 
\end{figure} 
 
\begin{figure} 
\centering\dontincludegraphics[angle=90,width=\linewidth]%
{figures/InspectZMContour_O3--v2-2_O3--v3-3_2006m11_LatPress-All-NaN-UT} 
\caption{Zonal averages for UTLS MLS ozone profiles during November,
  2006, showing the MLS v2.2 ozone mixing ratio contours (top left
  panel), the \thisv contours (top right panel), and their differences in
  ppmv (\thisv minus v2.2, bottom left panel) and percent (\thisv minus v2.2
  versus v2.2, bottom right panel). }
\label{fig:MLSv2v3O3ContourUTLS} 
\end{figure} 
 
Figures \ref{fig:MLSv2v3O3ContourStrat} and
\ref{fig:MLSv2v3O3ContourUTLS} show zonally averaged fields for the
(full) month of November, 2006, with only the properly screened profiles
from v2.2 and \thisv being used; mean differences (ppmv and percent) are
also shown.  The averages for stratospheric (and lower mesospheric)
ozone have typically not changed by more than 0.1\,ppmv,or 1 to 2\%; see
Figure \ref{fig:MLSv2v3O3ContourStrat}.  At low latitudes, zonal-mean
\thisv tends to be \texttildelow10\,ppbv larger than v2.2 from 215\,hPa to
100\,hPa, while at higher latitudes, the v2.2/\thisv relationship tends
to switch signs, with \thisv higher than v2.2 at 316\,hPa and 147\,hPa,
and lower at 215\,hPa and 100\,hPa (see Figure
\ref{fig:MLSv2v3O3ContourUTLS}).  Also, we note that the version 3.3
mean tropical profiles in the UTLS can exhibit significant systematic
vertical oscillations, mostly apparent during certain months (see the
`Artifacts' section below).

\paragraph*{Ozone Columns:} Changes in the MLS stratospheric ozone
columns (or in columns down to pressures between 100 and 316\,hPa) for
\thisv are quite small; typical daily zonal averages are within one
percent, and often within one DU.  There is no significant systematic
offset (offsets can change slightly between pressure levels and versus
latitude, including changes in sign).  The most significant difference
is the change in scatter, as observed for example in standard deviations
about zonal averages (with \thisv scatter often lower than in v2.2 by
several DU); this is largely a result of (and at the expense of) changes
in the horizontal smoothing, and poorer horizontal resolution for \thisv
data (see below).

 
\subsection*{Resolution} 
%\onestandardavk{O3_HR}{O\cs3} 
 
\begin{figure} 
\centering\dontincludegraphics[width=0.5\linewidth]{figures/avk-v3-3-O3_HR-35N} 
\caption{Typical two-dimensional (vertical and horizontal along-track)
  averaging kernels for the MLS \thisv O\cs3 at the equator; variation in
  the averaging kernels is sufficiently small that these are
  representative of typical profiles.  Colored lines show the averaging
  kernels as a function of MLS retrieval level, indicating the region of
  the atmosphere from which information is contributing to the
  measurements on the individual retrieval surfaces, which are denoted
  by plus signs in corresponding colors.  The dashed black line
  indicates the resolution, determined from the full width at half
  maximum (FWHM) of the averaging kernels, approximately scaled into
  kilometers (top axes).  (Upper) Vertical averaging kernels (integrated
  in the horizontal dimension for five along-track profiles) and
  resolution.  The solid black line shows the integrated area under each
  kernel (horizontally and vertically); values near unity imply that the
  majority of information for that MLS data point has come from the
  measurements, whereas lower values imply substantial contributions
  from a priori information.  (Lower) Horizontal averaging kernels
  (integrated in the vertical dimension) and resolution.  The averaging
  kernels are scaled such that a unit change is equivalent to one decade
  in pressure.}
\label{fig:avkO3Strat} 
\end{figure} 
 
\begin{figure} 
  \centering\dontincludegraphics[width=0.5\linewidth]{figures/avk-v3-3-O3_HR-35N-UTLS}
\caption{As for \protect{\ref{fig:avkO3Strat}} but zooming in on the
    upper troposphere and lower stratosphere region.}
\label{fig:avkO3UTLS} 
\end{figure} 
 
Vertical and horizontal smoothing constraints were changed for \thisv data
in an attempt to capitalize upon the higher vertical resolution offered
by the finer grid, while minimizing the vertical oscillations invited by
the additional vertical degrees of freedom.  Based on the width of the
averaging kernels shown in Figures~\ref{fig:avkO3Strat}
and~\ref{fig:avkO3UTLS}, the vertical resolution for the standard O\cs3
product is \texttildelow2.5\,km in the uppermost troposphere and stratosphere,
but degrades to 4 to 6\,km in the upper mesosphere and to \texttildelow3\,km at
316\,hPa. At best, lower stratospheric resolution can be about 2.3\,km,
which is an improvement over v2.2 data (but not by a factor of two --
the best resolution offered by the new 12-per-decade grid).  The
along-track resolution in the stratosphere has changed from \texttildelow
200\,km in v2.2 to 300 to 450\,km in \thisv, depending on altitude (the
upper stratosphere shows poorer resolution). In the mesosphere, this
along-track resolution varies between about 300 and 700\,km. In the
upper troposphere, the along-track resolution degrades from
\texttildelow300\,km at 120\,hPa to \texttildelow450\,km at 261\,hPa.  Typical
resolution values are provided in the summary Table~\ref{tab:O3Summary}.
The cross-track resolution is set by the 6\,km width of the MLS 240\,GHz
field of view.  The longitudinal separation of MLS measurements, set by
the Aura orbit, is 10\degsym--\,20\degsym\ over middle and lower
latitudes, with much finer sampling in polar regions.
 
\subsection*{Precision} 
 
The horizontal smoothing changes in \thisv, coupled with the finer
vertical retrieval grid, have led to some changes in estimated precision
(keeping in mind, however, the poorer \thisv horizontal resolution).  In
the upper stratosphere, the precisions are improved (smaller values) by
about 30\% in the \thisv data, although there is no need to quote
substantially different values than the (rounded off) values that are
given in Table~\ref{tab:O3Summary}.  In the UTLS, the Level 2 ozone
precision (uncertainty) values have worsened slightly (by \texttildelow20 to
30\%) from v2.2 to \thisv.  As found previously for v2.2 data, the Level~2
precision values are often slightly lower than the observed scatter in
the data, evaluated in a narrow latitude band centered around the
equator where atmospheric variability is expected to be small, or
obtained from a comparison between ascending and descending coincident
MLS profiles.
 
Negative precision values for ozone occur essentially for every data
point at pressures smaller than 0.01\,hPa, indicating increasing
influence from the \emph{a priori}, although some MLS information exists
(e.g., regarding average day/night differences) into the uppermost
mesosphere and lower thermosphere.  Generally, however, we recommend
that scientific studies be restricted to pressures of 0.02\,hPa or
larger.
 
\paragraph*{Column values:} As for the v2.2 retrievals, the estimated
precisions for the \thisv MLS column ozone abundances down to pressures of
100 to 215\,hPa are 2\% or less.  The typical empirical precision in the
columns based on (1-$\sigma$) variability in the tropics is 2 to 3\%.
However, see the comments above regarding the somewhat poorer
along-track (horizontal) resolution for \thisv data.
 
\subsection*{Accuracy} 
 
The accuracy estimates shown in Table~\ref{tab:O3Summary} are from an
analysis which propagated estimated systematic errors in MLS
calibration, spectroscopy, etc., through the v2.2 measurement system.
Results using the \thisv algorithms are not expected to differ
significantly.  The values shown here are intended to represent
2$\sigma$ estimates of accuracy.  Overall, we see no evidence, based on
a number of (published) comparisons with well established data sets,
that significant disagreements (outside the combined accuracy estimates)
or MLS-related issues exist for the MLS ozone product.  For more
details, see the MLS validation papers by \citet{FroidevauxEtal08O3},
\citet{JiangYEtal07}.  and \citet{LiveseyEtal08}, as well as references
therein; some more recent references relevant to MLS ozone and ozone
columns are available on the MLS website under `Publications'.  Future
validation studies using \thisv data will focus on longer-term changes and
on the UTLS region.
 
\paragraph*{Column values:} Sensitivity tests using systematic changes
in various parameters that could affect the accuracy of the MLS
retrievals lead to possible biases (2-$\sigma$ estimates) of about 4\%,
as an estimated accuracy for the MLS column values (from integrated MLS
ozone profiles down to 100, 147, and 215\,hPa).  See also the (v2.2)
validation papers (and subsequent ozone-related publications, e.g.,
available from the MLS website) for results on column ozone comparisons
versus satellite, sonde, and lidar data.
 
\subsection*{Data screening} 
\begin{description} 
 
\item[Pressure range:] \screeningrule{261\,--\,0.02\,hPa.}
  \standardrangestatement 
  \standardprecisionrule
  \evenstatusrule
\item[Quality:] \qualityrule{0.6}
\item[Convergence:] \convergencerule{1.18}
\item[Clouds:] \screeningrule{Scattering from thick clouds can lead to oscillatory
    values for O\cs3 in the UT/LS, see below for screening rules.}

  Most of the affected profiles are removed by the \texttt{Quality},
  \texttt{Convergence}, and outlier screening methods, as described below, 
  specifically for the \thisv data.
\item[`Outlier' Screening:] The \thisv ozone product has more tropical
  outliers than the v2.2 data in the lower stratosphere (mainly a few to
  a few tens of positive outliers at 68\,hPa on a typical day) and in
  the upper troposphere, where both negative or positive spikes can
  dominate, depending on the pressure level.  As for CO, these often
  appear to be related to thick clouds, and large values of ice water
  content (IWC) in the vicinity of the affected profiles.
  
  We have found that using explicit thresholds for negative ozone spikes
  at pressures larger than the 46\,hPa MLS pressure level can effectively
  screen out most of the outliers, including most of the \emph{positive}
  outliers in the UTLS.  The recommended screening method for \thisv ozone
  should therefore also include a test for significantly negative values
  for the MLS pressure range from 56 to 261\,hPa (inclusive).  We then
  recommend rejection of \emph{all values in this range} when a
  (negative) mixing ratio less than $-$0.15\,ppmv is encountered (for any
  level in this range); in addition, any value at 316\,hPa less than
  $-$0.30\,ppmv should also lead one to discard the applicable UTLS values.
  These outlier thresholds are chosen so that large negative values
  outside those thresholds are generally outside the `5 sigma level', in
  relation to typical ozone values and MLS ozone precisions (or scatter)
  in the UTLS; we have also checked that these criteria do not impact
  the screening of high latitude ozone values in any significant way
  (e.g., under ozone hole conditions).

  In summary, one should \emph{reject} profiles with odd \texttt{Status}
  \emph{or} even \texttt{Status} profiles with \texttt{Convergence}
  above the convergence threshold \emph{or} \texttt{Quality} below the
  quality threshold, or with values from 316 to 56\,hPa (inclusive) that
  get eliminated by the negative outlier criteria.  Conversely, one
  should \emph{keep} profile values with even status \emph{and} good
  \texttt{Convergence} \emph{and} good \texttt{Quality} \emph{and}, UTLS
  values in the 316 to 56\,hPa MLS pressure range that do not get
  discarded by the negative outlier criteria.  This methodology does
  remove some profile values that are not in the `outlier' category but
  no method will cleanly remove only the exact number of profiles that
  are suspicious for every day of the MLS mission.  The current
  recommendations typically remove \texttildelow4 to 6 \% of global daily data,
  with some tropical latitudes showing much larger fractional removal
  (e.g., 20 to 30\% in 10\degsym\ bins near the equator).  This
  screening generally maintains sufficient coverage for a near-complete
  daily map (for any given day), even in the UTLS\@.

  Compared to the v2.2 UT data screening recommendations, the screening
  of \thisv UTLS data generally removes slightly more ozone profiles on a
  typical day (although on occasion, slightly less).
 
\begin{figure} 
\centering\dontincludegraphics[angle=90,width=\linewidth]%
{figures/O3-vs-lat-unscreened-2005d343} 
\caption{Example of the MLS ozone mixing ratio distribution versus
  latitude for one day, showing various pressure levels in the UTLS and
  which points are flagged by recommended quality, convergence, and
  outlier criteria (for 2005d343, one of the days showing the most
  outliers). The bottom right panel (for 68\,hPa) also shows (red curve)
  the percentage of points getting screened out in 10\degsym\ latitude
  bins (with the y-axis scaled by a factor of 10) -- close to 30\% of
  points can be thrown out in some bins, on this relatively poor quality
  day.}
\label{fig:O3vsLatUnscreened} 
\end{figure} 
 
\begin{figure} 
\centering\dontincludegraphics[angle=90,width=\linewidth]%
{figures/O3-vs-lat-screened-2005d343} 
\caption{Example of the MLS ozone mixing ratio distribution versus
  latitude for one day, showing various pressure levels in the UTLS and
  which points remain \emph{after} removal of the flagged (colored)
  points shown in the previous Figure.}
\label{fig:O3vsLatScreened} 
\end{figure} 
 
 
Finally, we note that since there is essentially no impact from the
outliers (spikes) in the UTLS (this is also largely a tropical issue) on
the ozone mixing ratios at pressures less than 50\,hPa, it is safe to
ignore the outlier flag in the data screening if a study is only
concerned with pressures smaller than or equal to the 46\,hPa MLS
pressure level; in this case, users can simply apply the
\texttt{Quality} and \texttt{Convergence} (and \texttt{Status}) tests to
obtain a satisfactory data screening method, with fewer profiles removed
(typically less than about 2\% per day, globally) than if the UTLS
screening method were used.
\end{description} 
 
\subsection*{Review of comparisons with other datasets} 
 
O\cs3 comparisons have indicated general agreement at the 5\,--\,10\,\% 
level with stratospheric profiles from a number of comparisons using 
satellite, balloon, aircraft, and ground-based data.  A high MLS v2.2 
bias at 215\,hPa has been obtained in some comparisons versus 
ozonesondes, but this is not observed consistently in other comparisons.  
We have found that latitudinal and temporal changes observed in various 
correlative data sets are well reproduced by the MLS ozone product.  
Future (and ongoing) validation studies using \thisv data will focus more 
on longer-term changes and on the UTLS region, where improvements will 
still be sought, especially in the tropics (see below).
 
\subsection*{Artifacts} 

\begin{figure} 
  \centering\dontincludegraphics[width=\linewidth]%
  {figures/O3_mean_profiles_2006_April_July-v03-v02_25S-25N}
  \caption{Monthly zonal average MLS ozone profiles in the tropical UTLS
    (here, from 316 to 68\,hPa) for 2006 April (left panels) and July
    (right panels), for five 10\degsym \,latitude bins (top to bottom:
    15\degsym\,N--\,25\degsym\,N, 5\degsym\,N--\,15\degsym\,N,
    5\degsym\,S--\,5\degsym\,N, 15\degsym\,S--\,5\degsym\,S,
    25\degsym\,S--\,15\degsym\,S,).  The \thisv profiles are in red, while
    the v2.2 profiles are in blue.}
\label{fig:MLSO3UTtropicalOscillations} 
\end{figure} 

\begin{description}  
\item[Oscillations in tropical UTLS ozone:] The finer resolution and new
  retrieval methodology (and screening) for \thisv allow for improved
  values at 261\,hPa (and to some extent at 316\,hPa, although not a
  recommended level), but some artifacts (oscillations) exist in the
  tropical upper tropospheric profiles.  Indeed, vertical patterns exist
  in monthly means, as shown in Figure
  \ref{fig:MLSO3UTtropicalOscillations} with larger artifacts apparent
  in April (or May) than in July (or August), for example.  Further
  detailed validation and characterization of the MLS tropical UT data
  (in particular) is warranted, in order to more fully understand the
  MLS data limitations for various applications.
\item[Outliers:] Even with the data screening procedures that are
  recommended herein, a few outliers will remain unscreened for some
  days at some pressure levels (with the vast majority of outliers
  occurring at pressures larger than the MLS 46\,hPa level).  Caution is
  advised, not to over-interpret such occasional events.
\item[Columns:] Users of column ozone data above the tropopause from the
  MLS Level 2 files should be aware that the accuracy of these values
  depends on the tropopause pressure accuracy, and that artifacts can
  occur in these calculations, especially at high latitudes (under
  certain temperature gradient conditions).  Users should therefore
  inspect the MLS file values of tropopause pressure if using this
  product (swath) from the MLS ozone files.
\end{description}   
 
\subsection*{Priorities for future data version} 
\begin{itemize}  
\item Reduction of oscillations in UTLS tropical profiles, and retrieval
  improvements in the presence of thick clouds.
\end{itemize}
 
\begin{table} 
\caption{Summary for MLS ozone}\vspace{0.5ex} 
\label{tab:O3Summary} 
\begin{minipage}{\linewidth} 
\centering\begin{tabular}{ccccccl} 
\toprule 
\parbox[m]{15mm}{\centering\textbf{Pressure / hPa}} & 
\parbox[m]{24mm}{\centering\textbf{Resolution\\ Vert. $\times$ Horiz.}} & 
\multicolumn{2}{c}{\parbox[m]{20mm}{\centering\textbf{Precision\footnote{Precision  
        on individual profiles}\\ppmv\hfill\%}}} & 
\multicolumn{2}{c}{\parbox[m]{20mm}{\centering\textbf{Accuracy\footnote{%
        As estimated from systematic uncertainty characterization tests.
        Stratospheric values are expressed in ppmv with a typical
        \emph{equivalent} percentage value quoted.  215\,--\,100\,hPa errors are
        the sum of the ppmv \emph{and percentage} scaling uncertainties quoted.
        Accuracy values, especially for pressures from 100 to 316\,hPa will be
        re-evaluated, but the estimates for v2.2 data are currently used in this
        Table.}% 
      \\ppmv\hfill\%}}}  
& \textbf{Comments} \\ 
\midrule 
$\leq$ 0.01 & ---            & ---  & ---    & ---  & ---  & Unsuitable for scientific use \\   
0.02 & 5.5 $\times$ 200 & 1.4  & 300    & 0.1  &  35  & \\ 
0.05 & 5.5 $\times$ 200 & 0.9  & 150    & 0.2  &  30  & \\ 
0.1  & 4 $\times$ 400 & 0.5  &  60    & 0.2  &  20  & \\ 
0.2  & 3 $\times$ 450 & 0.5  &  40    & 0.1  &   7  & \\ 
0.5  & 3.5 $\times$ 550 & 0.3  &  20    & 0.1  &   5  & \\ 
1    & 3 $\times$ 450 & 0.2  &  7       & 0.2  &   7  & \\ 
2    & 3.5 $\times$ 450 & 0.15  &   3   & 0.2  &   5  & \\ 
5    & 3 $\times$ 450 & 0.15  &   2     & 0.3  &   5  & \\ 
10   & 3 $\times$ 500 & 0.1  &   2      & 0.3  &   5  & \\ 
22   & 2.5 $\times$ 400 & 0.1  &   2    & 0.2  &   5  & \\ 
46   & 2.5 $\times$ 350 & 0.06  &   3   & 0.2  &   8  & \\ 
68   & 2.5 $\times$ 350 & 0.04 & 3--10  & 0.05 & 3--10 & \\ 
100  & 2.5 $\times$ 300 & 0.04 & 20--30 & \multicolumn{2}{c}{$\left[0.05 +5\%\right]$} \\ 
150  & 2.5 $\times$ 400 & 0.03 & 5--100 & \multicolumn{2}{c}{$\left[0.02 +5\%\right]$} \\ 
215  & 3 $\times$ 400   & 0.02 & 5--100 & \multicolumn{2}{c}{$\left[0.02 +20\%\right]$}  \\ 
261  & 3 $\times$ 450   & 0.03 & 5--100 & ---  & ---   &  Requires further evaluation  \\ 
316  & 2.5 $\times$ 500 & 0.05 & ---    & ---  & ---   & Not recommended (until further evaluation) \\ 
1000\,--\,464 & ---     & ---  & ---    & ---  & ---   & Not retrieved \\ 
\bottomrule 
\end{tabular} 
\end{minipage} 
\end{table} 
 
%%% Local Variables:  
%%% mode: latex 
%%% TeX-master: "v4-x-report" 
%%% End:  
 
 
 
 
 
            
 
